\chapter{Ruang Vektor}

\Sage{} dapat bekerja dengan ruang vektor, misalnya dengan mencari basis
suatu ruang.

Dalam bab buku ini tentang ruang vektor, skalar berasal
dari himpunan bilangan real, $\Re$.
\Sage{} memungkinkan Anda memilih satu dari dua model bilangan real.
Yang pertama ialah \Sagecmd{RDF}, model bilangan real titik-mengambang
bawaan komputer.
% \footnote{%
%   The computer used to make this manual has
%   IEEE~754 double-precision binary floating-point numbers;  
%   if you have programmed then you may know this number model as binary64.}  
Yang lainnya ialah \Sagecmd{RR}, yang menyediakan bilangan real 
dengan presisi sebarang.
Model pertama berjalan lebih cepat
dan lebih luas digunakan dalam praktik, jadi itulah yang akan kita gunakan.
 




%========================================
\section{Ruang real}

\Sage{} memungkinkan Anda membuat ruang vektor.
Berikut ialah $\Re^4$.
\begin{sagecommandline}
sage: V = VectorSpace(RDF,4)
sage: V
\end{sagecommandline}
Anda juga dapat membuat subruang.
Berikut ialah subruang berdimensi satu dari $\Re^4$, yang dinyatakan sebagai rentang
himpunan dengan satu vektor.
\begin{sagecommandline}
sage: V = VectorSpace(RDF,4)
sage: V
sage: v1 = vector(RDF, [1, 1, -3, 0])
sage: W = V.span([v1])
sage: W
\end{sagecommandline}

Anda dapat mengajukan pertanyaan kepada \Sage{} tentang subruang ini, misalnya 
menguji keanggotaan. 
\begin{sagecommandline}
sage: v3 = vector(RDF, [2, 2, -6, 0])
sage: v3 in W
sage: v4 = vector(RDF, [1, 0, 0, 5])
sage: v4 in W
\end{sagecommandline}

Contoh lain ialah~$\Re^3$.
\begin{sagecommandline}
sage: V = VectorSpace(RDF,3)
sage: V
sage: v1 = vector(RDF, [1, 2, 3])
sage: v1 in V
sage: v2 = vector(RDF, [1, 2, 3, 4])
sage: v2 in V
\end{sagecommandline}

Kita dapat membuat subruang sebagai rentang beberapa vektor.
\begin{sagecommandline}
sage: V = VectorSpace(RDF,3)
sage: v1 = vector(RDF, [1, 0, 3])
sage: v2 = vector(RDF, [0, 1, 1])
sage: W = V.span([v1, v2])
sage: W
sage: v3 = vector(RDF, [1, 1, 4])
sage: v3 in W
\end{sagecommandline}



\subsection{Basis}
\Sage{} akan memperoleh basis suatu ruang vektor atau subruang.
\begin{sagecommandline}
sage: V = VectorSpace(RDF,3)
sage: V.basis()
sage: v1 = vector(RDF, [2, -1, -3])
sage: v2 = vector(RDF, [1, 1, 0])
sage: W = V.span([v1, v2])     
sage: W.basis()
\end{sagecommandline}
Perhatikan bahwa basis yang dikembalikan oleh \Sage{}  
tidak sama dengan himpunan dua vektor
$\set{\vec{v}_1, \vec{v}_2}$ yang kita berikan kepadanya.
Untuk memperoleh basis baru, \Sage{} menempatkan vektor-vektor dari himpunan perentang 
sebagai baris-baris suatu matriks,
mengubah matriks itu ke bentuk eselon baris tereduksi, lalu melaporkan baris-baris 
taknol sebagai basis.
Setiap matriks memiliki tepat satu bentuk eselon baris tereduksi, sehingga setiap 
subruang dari~$V$ memiliki tepat satu basis semacam itu;
inilah basis kanonis subruang tersebut.

\Sage{} menampilkan basis ini ketika Anda meminta deskripsi
subruang tersebut.
\begin{sagecommandline}
sage: W  
\end{sagecommandline}

Subruang lain dari $\Re^3$ ialah bidang yang dinyatakan oleh persamaan
$x-2y+2z=0$.
\begin{equation*}
  W=\set{\colvec{x \\ y \\ z}
    \suchthat x=2y-2z}
  =\set{\colvec{2 \\ 1 \\ 0}y+\colvec[r]{-2 \\ 0 \\ 1}z
        \suchthat y,z\in\Re}
\end{equation*}
Berikut subruang tersebut.
\begin{sagecommandline}
sage: V = VectorSpace(RDF,3)               
sage: v1 = vector(RDF, [2, 1, 0]) 
sage: v2 = vector(RDF, [-2, 0, 1]) 
sage: W = V.span([v1, v2])       
sage: W.basis()
\end{sagecommandline}

Jika Anda mengambil himpunan~$\set{\vec{v}_1,\vec{v}_2}$ itu dan menambahkan vektor~$\vec{v}_3$ yang
bebas linear dari kedua vektor tersebut, 
himpunan itu menghasilkan subruang berdimensi tiga dari $\Re^3$, 
yaitu seluruh $\Re^3$.
\begin{sagecommandline}
sage: v3 = vector(RDF, [2, 3, 0])
sage: W_prime = V.span([v1, v2, v3])
sage: W_prime.basis()
\end{sagecommandline}
Sebaliknya, jika Anda menambahkan ke himpunan perentang $\set{\vec{v}_1,\vec{v}_2}$ 
vektor ketiga~$\vec{v}_3$ 
yang bergantung linear pada kedua vektor lainnya, maka
rentangnya tidak berubah.
Artinya, dalam kasus ini subruang yang direntang oleh dua vektor pertama,
$\spanof{\vec{v}_1,\vec{v}_2}$,
sama dengan rentang ketiganya,
$\spanof{\vec{v}_1,\vec{v}_2,\vec{v}_3}$.
\begin{sagecommandline}
sage: W = V.span([v1, v2])       
sage: W.basis()
sage: v3 = vector(RDF, [0, 1, 1])
sage: W_prime = V.span([v1, v2, v3])
sage: W_prime.basis()
\end{sagecommandline}

Jika Anda ingin menggunakan basis sendiri, \Sage{} juga
memungkinkan hal itu.
\begin{sagecommandline}
sage: V = VectorSpace(RDF,3)
sage: v1 = vector(RDF, [1, 2, 3])
sage: v2 = vector(RDF, [2, 1, 3])
sage: W = V.span_of_basis([v1, v2])
sage: W.basis()
\end{sagecommandline}




\subsection{Kesamaan}

\Sage{} dapat membandingkan ruang.
Berikut dua deskripsi bidang-$xy$ di dalam $\Re^4$.
\begin{sagecommandline}
sage: V = VectorSpace(RDF,4)
sage: v1 = vector(RDF, [1, 0, 0, 0])
sage: v2 = vector(RDF, [1, 1, 0, 0])
sage: W12 = V.span([v1, v2])
sage: v3 = vector(RDF, [2, 1, 0, 0])
sage: W13 = V.span([v1, v3])  
\end{sagecommandline}
% We could test that the two spaces $W_{1,2}$ and $W_{1,3}$ 
% are equal by using set membership.
Jika kedua vektor yang digunakan untuk membentuk himpunan perentang $W_{1,2}$ 
merupakan anggota $W_{1,3}$, maka $W_{1,2}\subseteq W_{1,3}$.
Jika, selain itu, vektor-vektor yang digunakan untuk membentuk himpunan perentang $W_{1,3}$ 
merupakan anggota $W_{1,2}$, maka kedua ruang itu sama.
\begin{sagecommandline}
sage: v2 in W13
sage: v3 in W12
\end{sagecommandline}
Karena $\vec{v}_1,\vec{v}_3\in W_{1,2}$ dan 
$\vec{v}_1,\vec{v}_2\in W_{1,3}$, kedua subruang itu sama.

Namun, cara yang lebih langsung untuk menguji kesamaan ialah dengan menanyakannya saja.
\begin{sagecommandline}
sage: W12 == W13
\end{sagecommandline}
Penggunaan operator kesamaan, \inlinecode{==}, ini 
belum lengkap tanpa mencobanya pada dua ruang yang
tidak sama. 
\begin{sagecommandline}
sage: v4 = vector(RDF, [1, 1, 1, 1])
sage: W14 = V.span([v1, v4])
sage: v2 in W14
sage: v4 in W12
sage: W12 == W14
sage: W12 != W14
\end{sagecommandline}

Ada hal penting tentang algoritme di sini.
\Sage{} dapat memeriksa kesamaan dua rentang dengan sejumlah cara.
Salah satunya ialah memeriksa apakah setiap anggota himpunan perentang pertama berada di
ruang kedua, dan apakah setiap anggota himpunan perentang kedua berada di 
ruang pertama. 
Namun, \Sage{} melakukan hal yang berbeda.
Ia memeriksa kesamaan ruang
hanya dengan memeriksa apakah basis kanonisnya sama.
\begin{sagecommandline}
sage: W12.basis()
sage: W14.basis()
\end{sagecommandline}
Kedua algoritme ini 
memiliki perilaku eksternal yang sama, karena keduanya memutuskan apakah
dua ruang sama.
Namun, kedua algoritme itu berbeda secara internal dan yang kedua lebih cepat, 
sebagian karena \Sage{} dapat menghitung basis kanonis suatu ruang terlebih dahulu
lalu menggunakannya sebanyak yang diperlukan.
Menemukan cara tercepat untuk menyelesaikan tugas merupakan bidang penelitian penting dalam komputasi.


\subsection{Operasi}
\Sage{} dapat mencari irisan dua ruang.
Perhatikan anggota-anggota $\Re^3$ berikut.
\begin{equation*}
  \vec{v}_1=\colvec{1 \\ 0 \\ 0}
  \quad \vec{v}_2=\colvec{0 \\ 1 \\ 0}
  \quad \vec{v}_3=\colvec{0 \\ 0 \\ 3}
\end{equation*}
Bentuklah dua rentang, bidang-$xy$ $W_{1,2}=\spanof{\vec{v}_1, \vec{v}_2}$ 
dan bidang-$yz$ $W_{2,3}=\spanof{\vec{v}_2, \vec{v}_3}$.
Irisan keduanya ialah sumbu-$y$. 
\begin{sagecommandline}
sage: V = VectorSpace(RDF,3)
sage: v1 = vector(RDF, [1, 0, 0])
sage: v2 =  vector(RDF, [0, 1, 0])
sage: W12 = V.span([v1, v2])
sage: W12.basis()
sage: v3 = vector(RDF, [0, 0, 3])
sage: W23 = V.span([v2, v3])
sage: W23.basis()
sage: W = W12.intersection(W23)
sage: W.basis()
\end{sagecommandline}

Jika Anda mencoba mengiris dua ruang ketika operasi itu tidak bermakna,
misalnya mencoba mengiris $\Re^3$ dan~$\Re^4$, maka \Sage{}
memberikan pesan galat yang baris terakhirnya memuat string
\inlinecode{self and other must have the same ambient space.}

Ingat bahwa ruang trivial $\set{\zero}$ merupakan rentang himpunan kosong.
\begin{sagecommandline}
sage: V = VectorSpace(RDF,3)
sage: v1, v2 = vector(RDF, [1, 0, 0]), vector(RDF, [0, 1, 0])
sage: W12 = V.span([v1, v2])
sage: v3 = vector(RDF, [1, 1, 1])
sage: W3 = V.span([v3])
sage: W3.basis()
sage: W4 = W12.intersection(W3)
sage: W4.basis()
\end{sagecommandline}

\Sage{} juga akan mencari jumlah ruang, yaitu rentang gabungannya.
\begin{sagecommandline}
sage: V = VectorSpace(RDF,3)
sage: v1, v2 = vector(RDF, [1, 0, 0]), vector(RDF, [0, 1, 0])
sage: W12 = V.span([v1, v2])
sage: v3 = vector(RDF, [1, 1, 1])
sage: W3 = V.span([v3])
sage: W5 = W12 + W3
sage: W5
sage: W5 == V
\end{sagecommandline}
(Buku membahas jumlah subruang dalam sebuah bagian opsional.)





%========================================
\section{Ruang lain}

Di dalam buku terdapat ruang-ruang vektor yang bukan subruang dari suatu $\Re^n$.
Untuk bekerja dengannya, Anda dapat menggunakan sejumlah cara yang lebih canggih, tetapi
salah satu cara yang langsung ialah 
menggunakan ruang real yang serupa dengan ruang yang diinginkan.\footnote{%
  Bab ketiga buku teks ini, 
  Pemetaan Antar-Ruang, memberikan makna yang tepat bagi 
  ``serupa''.}
Perhatikan subruang dari 
$\polyspace_2$.
\begin{equation*}
  \set{ a_2x^2+a_1x+a_0\suchthat a_2=a_0+a_1}           
   =\set{ (a_1+a_0)x^2+a_1x+a_0 \suchthat a_1,a_0\in\Re}
\end{equation*}
Ruang itu serupa dengan
subruang dari $\Re^3$ berikut.
\begin{equation*}
  \set{\colvec{a_1+a_0 \\ a_1 \\ a_0}\suchthat a_1,a_0\in\Re}
  =\set{\colvec{1 \\ 1 \\ 0}a_1+\colvec{1 \\ 0 \\ 1}a_0\suchthat a_1,a_0\in\Re}
\end{equation*}
\begin{sagecommandline}
sage: V = VectorSpace(RDF,3)
sage: v1 = vector(RDF, [1, 1, 0])
sage: v2 = vector(RDF, [1, 0, 1])
sage: W = V.span([v1, v2])
sage: W.basis()
\end{sagecommandline}

Dengan cara serupa, Anda dapat menggunakan \Sage{} untuk 
melakukan perhitungan pada ruang matriks $\nbyn{2}$ berikut,
\begin{equation*}
  \set{\begin{mat}
         a  &b \\
         c  &d
       \end{mat}\suchthat \text{$a-b+c=0$ dan $b+d=0$}}
\end{equation*}
dengan melakukan parameterisasi
\begin{equation*}
  \set{\begin{mat}
         a  &b \\
         c  &d
       \end{mat}\suchthat \text{$a=-c-d$ dan $b=-d$}}
  =\set{\begin{mat}
         -1  &0 \\
          1  &0
       \end{mat}c
       +
       \begin{mat}
         -1  &-1 \\
          0  &1
       \end{mat}d
       \suchthat c,d\in\Re}
\end{equation*}
lalu memberikan ruang real yang bersesuaian kepada \Sage{}.
\begin{sagecommandline}
sage: V = VectorSpace(RDF,4)
sage: v1 = vector(RDF, [-1, 0, 1, 0])
sage: v2 = vector(RDF, [-1, -1, 0, 1])
sage: W = V.span([v1, v2])
sage: W.basis()
\end{sagecommandline}
Contoh itu memperoleh vektor dengan membaca sepanjang baris-baris matriks.
Sebagai gantinya, Anda dapat membaca
menuruni kolom-kolomnya.  
\begin{sagecommandline}
sage: V = VectorSpace(RDF,4)
sage: v1 = vector(RDF, [-1, 1, 0, 0])
sage: v2 = vector(RDF, [-1, 0, -1, 1])
sage: W = V.span([v1, v2])
sage: W.basis()
\end{sagecommandline}
Masih ada cara-cara lain untuk menghasilkan ruang yang bersesuaian.
Ruang-ruang itu tampak berbeda satu sama lain,
tetapi sifat-sifat pentingnya, seperti dimensi, 
tidak dipengaruhi oleh tampilannya.
Pembahasan yang jauh lebih lengkap tentang hal ini terdapat dalam bab ketiga buku.

\endinput


TODO:
